\chapter{Hauptmengen, Unter- und Obermengen}
\label{b12guide:ch:order-sets}

\section{Hauptmengen und Wiedergewinnung der Ordnung}

\begin{definition}[Untere und obere Hülle; Hauptmengen]
\label{b12guide:def:closures}
Für $S\subseteq P$ seien
\[
\begin{aligned}
 \downarrow S&=\{x\in P\mid\exists s\in S\;(x\leq s)\},\\
 \uparrow S&=\{x\in P\mid\exists s\in S\;(s\leq x)\}.
\end{aligned}
\]
Für $a\in P$ schreiben wir
$\downarrow a=\downarrow\{a\}$ und
$\uparrow a=\uparrow\{a\}$. Diese Mengen heißen
\emph{Hauptuntermenge} und \emph{Hauptobermenge} von $a$;
$\uparrow a$ ist der \emph{Hauptfilter} aus Band~12.
Bei mehreren Ordnungen wird der Träger oder der Relationsgraph
als Index angegeben.
Die Definitionen verwenden einen Existenzquantor; sie bezeichnen
keine Mengen gemeinsamer Schranken von $S$.
\end{definition}

\begin{theorem}[Extremalpunkte der Hauptmengen]
\label{b12guide:thm:principal-extremes}
Für $a\in P$ gilt
\[
 \max(\downarrow a)=a,\qquad
 \min(\uparrow a)=a,\qquad
 (\downarrow a)\cap(\uparrow a)=\{a\}.
\]
Insbesondere ist $a$ das einzige maximale Element von
$\downarrow a$ und das einzige minimale Element von $\uparrow a$.
\end{theorem}

\begin{theorem}[Ordnungs- und Gleichheitskriterien durch Hauptmengen]
\label{b12guide:thm:principal-recovery}
Für $x,y\in P$ gilt
\[
\begin{aligned}
 x\leq y
 &\quad\Longleftrightarrow\quad\downarrow x\subseteq\downarrow y
 \quad\Longleftrightarrow\quad\uparrow y\subseteq\uparrow x,\\
 x=y
 &\quad\Longleftrightarrow\quad\downarrow x=\downarrow y
 \quad\Longleftrightarrow\quad\uparrow x=\uparrow y,\\
 x<y
 &\quad\Longleftrightarrow\quad\downarrow x\subsetneq\downarrow y
 \quad\Longleftrightarrow\quad\uparrow y\subsetneq\uparrow x.
\end{aligned}
\]
Damit ist jeder Vergleich vollständig durch Inklusionen zwischen
Hauptmengen bestimmt.
\end{theorem}

\section{Abgeschlossenheit nach unten und nach oben}

\begin{definition}[Untere und obere Teilmenge]
\label{b12guide:def:lower-upper}
Eine Teilmenge $L\subseteq P$ heißt \emph{untere Menge}, wenn
\[
 x\in L,\ y\in P,\ y\leq x\quad\Longrightarrow\quad y\in L.
\]
Eine Teilmenge $U\subseteq P$ heißt \emph{obere Menge}, wenn
\[
 x\in U,\ y\in P,\ x\leq y\quad\Longrightarrow\quad y\in U.
\]
Die Familien aller unteren und aller oberen Mengen heißen
$\Low(P)$ und $\Up(P)$. Es wird keine zusätzliche
Gerichtetheit oder Nichtleerheit verlangt.
\end{definition}

\begin{theorem}[Kleinste untere und obere Menge]
\label{b12guide:thm:least-closure}
Für $S\subseteq P$ ist $\downarrow S$ die kleinste untere Menge,
die $S$ enthält, und $\uparrow S$ die kleinste obere Menge,
die $S$ enthält. Für jede untere Menge $L$ und jede obere Menge $U$ gilt also
\[
 S\subseteq L\ \Longleftrightarrow\ \downarrow S\subseteq L,
 \qquad
 S\subseteq U\ \Longleftrightarrow\ \uparrow S\subseteq U.
\]
\end{theorem}

\begin{theorem}[Charakterisierung durch Hauptmengen und Hüllen]
\label{b12guide:thm:closed-criteria}
Für $S\subseteq P$ gilt
\[
\begin{aligned}
 S\in\Low(P)
 &\quad\Longleftrightarrow\quad\downarrow S=S
 \quad\Longleftrightarrow\quad
       \forall s\in S\;(\downarrow s\subseteq S),\\
 S\in\Up(P)
 &\quad\Longleftrightarrow\quad\uparrow S=S
 \quad\Longleftrightarrow\quad
       \forall s\in S\;(\uparrow s\subseteq S).
\end{aligned}
\]
Jede untere Menge $L$ und jede obere Menge $U$ besitzen die
Darstellungen
\[
 L=\bigcup_{x\in L}\downarrow x,\qquad
 U=\bigcup_{x\in U}\uparrow x.
\]
\end{theorem}

\begin{theorem}[Hüllengesetze]\label{b12guide:thm:closure-laws}
Für $S,T\subseteq P$ gilt
\[
\begin{aligned}
 S\subseteq T&\ \Rightarrow\
 \downarrow S\subseteq\downarrow T\ \land\
 \uparrow S\subseteq\uparrow T,\\
 \downarrow(\downarrow S)&=\downarrow S,\qquad
 \uparrow(\uparrow S)=\uparrow S.
\end{aligned}
\]
Für jede Familie $(S_i)_{i\in I}$ von Teilmengen von $P$ gilt
\[
 \downarrow\!\left(\bigcup_{i\in I}S_i\right)
     =\bigcup_{i\in I}\downarrow S_i,\qquad
 \uparrow\!\left(\bigcup_{i\in I}S_i\right)
     =\bigcup_{i\in I}\uparrow S_i.
\]
Dies schließt $I=\varnothing$ ein:
$\downarrow\varnothing=\uparrow\varnothing=\varnothing$.
Weiter gilt $\downarrow P=\uparrow P=P$.
\end{theorem}

\begin{theorem}[Vereinigung, Durchschnitt und Komplement]
\label{b12guide:thm:set-operations}
Beliebige Vereinigungen und beliebige Durchschnitte unterer Mengen
sind untere Mengen. Dasselbe gilt für obere Mengen.
Der leere Durchschnitt wird dabei relativ zu $P$ als $P$
verstanden; die leere Vereinigung ist $\varnothing$.
Außerdem gilt für jedes $S\subseteq P$
\[
 S\in\Low(P)\quad\Longleftrightarrow\quad P\setminus S\in\Up(P).
\]
Beim Dualisieren werden die beiden Familien und beide Hüllen
vertauscht:
\[
 \Low(P^{\mathrm{op}})=\Up(P),\qquad
 \downarrow_{P^{\mathrm{op}}}S=\uparrow_P S.
\]
\end{theorem}

\section{Darstellung der Ordnung und Extremalmengen}

\begin{theorem}[Trennung und Rekonstruktion durch untere und obere Mengen]
\label{b12guide:thm:separation}
Für $x,y\in P$ sind äquivalent:
\[
\begin{aligned}
 &x\leq y,\\
 &\forall U\in\Up(P)\;(x\in U\Rightarrow y\in U),\\
 &\forall L\in\Low(P)\;(y\in L\Rightarrow x\in L).
\end{aligned}
\]
Ist $x\nleq y$, so gibt es eine obere Menge, die $x$, aber nicht
$y$ enthält, und eine untere Menge, die $y$, aber nicht $x$ enthält.
\end{theorem}

\begin{theorem}[Vergleich zweier Ordnungen durch ihre Teilmengen]
\label{b12guide:thm:refinement}
Für $R,S\in\PO(P)$ sind äquivalent:
\[
\begin{aligned}
 &R\subseteq S,\\
 &\forall x\in P\;(\downarrow_R x\subseteq\downarrow_S x),\\
 &\forall x\in P\;(\uparrow_R x\subseteq\uparrow_S x),\\
 &\Low(P,\leq_S)\subseteq\Low(P,\leq_R),\\
 &\Up(P,\leq_S)\subseteq\Up(P,\leq_R).
\end{aligned}
\]
Insbesondere bestimmen die Familie aller unteren Mengen und die
Familie aller oberen Mengen jeweils die Halbordnung eindeutig.
\end{theorem}

\begin{theorem}[Extremalelemente der erzeugten Mengen]
\label{b12guide:thm:closure-extremal}
Für jede Teilmenge $S\subseteq P$ gilt
\[
 \MinP(\uparrow S)=\MinP(S),\qquad
 \MaxP(\downarrow S)=\MaxP(S).
\]
Ist $C$ eine Antikette, so gilt insbesondere
\[
 \MinP(\uparrow C)=C,\qquad \MaxP(\downarrow C)=C.
\]
Für Antiketten $C,D$ folgt sowohl aus
$\uparrow C=\uparrow D$ als auch aus
$\downarrow C=\downarrow D$ bereits $C=D$.
\end{theorem}

\begin{theorem}[Wann minimale Elemente eine obere Menge erzeugen]
\label{b12guide:thm:minimal-generation}
Für $U\in\Up(P)$ gilt $U=\uparrow\MinP(U)$ genau dann,
wenn jedes $u\in U$ oberhalb mindestens eines minimalen Elements
von $U$ liegt. Dual gilt für $L\in\Low(P)$ die Gleichheit
$L=\downarrow\MaxP(L)$ genau dann, wenn jedes $x\in L$
unterhalb mindestens eines maximalen Elements von $L$ liegt.
Diese Bedingungen sind nicht automatisch erfüllt.
\end{theorem}
